\chapter{Fragmentarea relațiilor}
\label{chap:fragmentare}

\section*{a. Fragmentare orizontală primară}

Fragmentarea orizontală primară a fost aplicată pe relația \textit{CLIENT}, având ca scop separarea clienților în funcție de tipul acestora pentru a asigura localizarea datelor pe nodurile specifice tipului de business.

\textbf{i. Aplicarea algoritmului de fragmentare orizontală primară}\\
Algoritmul pornește de la analiza cererilor predominante din sistem:
\begin{itemize}
    \item $q_1$: Aplicația portalului B2C, care procesează datele persoanelor fizice.
    \item $q_2$: Aplicația portalului B2B, care procesează datele companiilor.
\end{itemize}

Fie atributul de partiționare \textit{client\_type}. Definim setul de predicate simple $Pr = \{p_1, p_2\}$:
\begin{itemize}
    \item $p_1 : \text{client\_type} = \text{'F'}$
    \item $p_2 : \text{client\_type} = \text{'J'}$
\end{itemize}

Setul $Pr$ este complet și minimal deoarece acoperă toate tipurile posibile de clienți din sistem și nu conține predicate redundante. Din acest set, generăm setul de predicate minterm $M = \{m_1, m_2\}$ prin combinarea predicatelor simple:
\begin{itemize}
    \item $m_1 = p_1 \wedge \neg p_2 \Rightarrow \text{client\_type} = \text{'F'}$ (Minterm Valid)
    \item $m_2 = \neg p_1 \wedge p_2 \Rightarrow \text{client\_type} = \text{'J'}$ (Minterm Valid)
\end{itemize}

Analizând frecvențele de acces, observăm că aplicația $q_1$ accesează exclusiv datele ce satisfac mintermul $m_1$, iar aplicația $q_2$ accesează exclusiv datele ce satisfac mintermul $m_2$. Mintermii care ar genera intersecții vide sau care nu sunt accesați de nicio aplicație au fost eliminați.

\textbf{ii. Obținerea fragmentelor orizontale primare}\\
Prin aplicarea operatorului de selecție ($\sigma$) pe baza mintermilor validați, obținem următoarele fragmente poziționate pe nodurile corespunzătoare:
\begin{itemize}
    \item $CLIENT\_FIZIC = \sigma_{\text{client\_type}='F'}(CLIENT)$ --- Alocat pe \textbf{Serverul 1 (SV1)}.
    \item $CLIENT\_JURIDIC = \sigma_{\text{client\_type}='J'}(CLIENT)$ --- Alocat pe \textbf{Serverul 2 (SV2)}.
\end{itemize}

\section*{b. Fragmentare orizontală derivată}

\textbf{i. Obținerea fragmentelor orizontale derivate}\\
S-a aplicat fragmentarea orizontală derivată pe relația \textit{TICKET}, deoarece tichetele nu pot fi partiționate direct printr-un atribut propriu fără a rupe legătura de coerență cu tipul de client. Fragmentarea se obține aplicând operația de semi-join ($\ltimes$) între relația copil (\textit{TICKET}) și fragmentele primare ale relației părinte (\textit{CLIENT}):
\begin{itemize}
    \item $TICKET\_FIZIC = TICKET \ltimes_{\text{client\_id}} CLIENT\_FIZIC$ --- Alocat pe \textbf{Serverul 1 (SV1)}.
    \item $TICKET\_JURIDIC = TICKET \ltimes_{\text{client\_id}} CLIENT\_JURIDIC$ --- Alocat pe \textbf{Serverul 2 (SV2)}.
\end{itemize}
Prin această abordare, tichetele deschise de persoanele fizice sunt stocate local pe același nod cu datele de profil ale clienților respectivi, optimizând masiv operațiile de join inter-noduri.

\section*{c. Fragmentare verticală}

Fragmentarea verticală a fost aplicată pe relația \textit{AGENT}. Din motive stringente de securitate și izolare a datelor, s-a decis separarea credentialelor sensibile de autentificare de datele publice de profil utilizate curent în aplicație.

\textbf{i. Aplicarea algoritmului de fragmentare verticală}\\
Relația globală \textit{AGENT} conține schema de atribute: $A_1$ (\textit{agent\_id}), $A_2$ (\textit{nume}), $A_3$ (\textit{prenume}), $A_4$ (\textit{telefon}), $A_5$ (\textit{email}), $A_6$ (\textit{password\_hash}).

Considerăm aplicațiile/modulele care accesează aceste atribute și frecvențele lor ipotetice de acces:
\begin{itemize}
    \item $q_1$ (Modulul Operațional/Afișare profil): Accesează în mod repetat setul de atribute $U_1 = \{A_1, A_2, A_3, A_4\}$.
    \item $q_2$ (Modulul de Securitate/Autentificare): Accesează strict la login setul $U_2 = \{A_1, A_5, A_6\}$.
\end{itemize}

Pe baza aplicațiilor, se construiește \textit{Matricea de Afinitate a Atributelor} (ipotetică), unde afinitatea $aff(A_i, A_j)$ măsoară legătura dintre atribute. Atributele $\{A_2, A_3, A_4\}$ prezintă o afinitate strânsă între ele datorită interogărilor frecvente din $q_1$, în timp ce $\{A_5, A_6\}$ sunt intens legate prin interogările de login din $q_2$. 

Aplicând algoritmul Bond Energy (BEA) pentru gruparea atributelor cu afinitate mare, se generează două clustere distincte de atribute disjuncte, ambele păstrând cheia primară $A_1$ pentru a asigura proprietatea de reconstrucție.

\textbf{ii. Obținerea fragmentelor verticale}\\
Prin aplicarea operatorului de proiecție ($\pi$) pe baza clusterelor identificate de algoritm, obținem următoarele două fragmente structurale:
\begin{itemize}
    \item $AGENT\_PROFIL = \pi_{\text{agent\_id, nume, prenume, telefon}}(AGENT)$ --- Conține datele publice de contact ale agenților. Aceste date sunt accesibile de pe nodurile operaționale (\textbf{SV1} și \textbf{SV2}).
    \item $AGENT\_SEC = \pi_{\text{agent\_id, email, password\_hash}}(AGENT)$ --- Conține exclusiv datele critice de acces. Acest fragment este izolat complet pe \textbf{Serverul 3 (SV3 - Modulul de Securitate)}.
\end{itemize}